\documentclass[a4paper,12pt]{article}
\usepackage[utf8]{inputenc}
\usepackage[T1]{fontenc}
\usepackage[french]{babel}
\usepackage{fancyhdr}
\usepackage{hyperref}
\usepackage{graphicx}
\usepackage{float}
\usepackage{times}

\title{Rapport}
\author{Alexandre Lou-Poueyou \\
Aymeric Cassard\\
Ludovic Andrieu\\
Jean Mayol\\
Matheo Petry\\ 
Valentin Cazet}

\date{\today}

\begin{document}

\maketitle

%----------------------------------------------------------------------------------------------------------------

\begin{abstract}
Ce document est le rapport du projet Scrabble en Java. 
\end{abstract}

\tableofcontents

%----------------------------------------------------------------------------------------------------------------

\section{Le contexte et l'existant du projet}

%----------------------------------------------------------------------------------------------------------------

\section{Une explication détaillée du sujet}

\subsection{Règles du jeu}

Le Scrabble est un jeu de société qui consiste à former des mots sur
une grille carrée de 15x15 cases.
Les 100 lettres utilisées ont une valeur de point choisie selon la
fréquence du mot dans la langue.
2 lettres 'Joker' supplémentaires sont présentes, elles permettent
au joueur de placer n'importe quelle lettre, mais n'ajoutent aucun
point au mot.
Chaque joueur dispose de lettres tirées aléatoirement qu'il place
sur la grille pour marquer des points.\newline
La grille contient des cases bonus:
\begin{itemize}
  \item Lettre double: double la valeur de la lettre posée.
  \item Lettre triple: triple la valeur de la lettre posée.
  \item Mot double: double la valeur du mot posé.
  \item Mot triple: triple la valeur du mot posé.
\end{itemize}

Les multiplicateurs de mots se combinent dans le cas où le mot posé
touche plusieurs cases.

Le mot doit exister dans le dictionnaire courant du programme pour
pouvoir être posé.
Les éventuels mots adjacents au mot posé rapportent des points.

Chaque joueur dispose de 7 lettres dans sa main.
Si un joueur pose toutes ses lettres en un tour, alors s'applique
un bonus de 50 points.\newline
Les actions que peut entreprendre le joueur lors de son tour sont:
\begin{itemize}
  \item Poser un mot.
  \item Échanger des lettres de la pioche.
  \item Passer son tour.
\end{itemize}

Les conditions de fin de partie sont les suivantes:
\begin{itemize}
  \item \textbf{Reliquat épuisé}: Un joueur a épuisé toutes ses
        lettres et la pioche est vide.
  \item \textbf{Blocage}: Les joueurs ne trouvent plus de mots a
        jouer et il reste moins de 7 lettres dans la pioche.
        Les valeurs des lettres de leur mains sont déduites de leur
        score final.
  \item \textbf{Temps épuisé}: Les deux joueurs n'ont plus de temps
        (variante Blitz uniquement).
\end{itemize}

Le joueur avec le score le plus élevé est le gagnant.\newline\newline
\textit{Variantes}
\begin{itemize}
  \item \textbf{Super-scrabble}: La grille fait 21x21.
  \item \textbf{Blitz}: Chaque joueur dispose d'un temps imparti, qui
        s'écoule lors de son tour.
        Si un joueur n'a plus de temps, il ne peut plus jouer.
\end{itemize}

\subsection{Structures de données}

\textit{Structures de données simples}\newline
L'utilisation de Java dans le projet permet d'éviter 
d'implémenter par nous-mêmes les structures de données de base 
du projet.
Par exemple, la pioche est représentée par une HashTable qui
associe chaque lettre à un compteur, représentant le nombre de
lettres restantes dans la pioche.\newline

\textit{Représentation du dictionnaire}\newline
Une manière commune de représentation de lexique est un arbre pour 
lequel chaque arête est une lettre du mot, et certains nœuds sont
marqués comme finaux, indiquant que le chemin jusqu'à ce noeud
constitue un mot. Cette structure est communément appelée un trie.
\begin{figure}[H]
    \centering
    \includegraphics[scale=.5]{img/trie.png}
\end{figure}
Le trie n'est pas utilisé dans notre projet mais aide à visualiser
le fonctionnement des structures de données suivantes.\newline

Le \textbf{DAWG}(Directed Acyclic Word-Graph) nous permet de
représenter le dictionnaire de validité des mots en mémoire.
Il s'agit d'un automate fini pour lequel le langage est l'ensemble
des mots du dictionnaire, et les transitions sont les lettres que
l'on ajoute pour former le mot. Enfin, les états finaux sont les
différents mots du lexique.
\begin{figure}[H]
    \centering
    \includegraphics[scale=.5]{img/dawg.png}
\end{figure}
Cette structure de données est équivalente au trie en termes  
de fonctionnalités, mais prend moins de place en mémoire.
\newline

Le \textbf{GADDAG} est une autre manière de représenter le lexique.
Dans cette structure, chaque mot est représenté à partir de chacune
de ses lettres intermédiaires.
\begin{figure}[H]
    \centering
    \includegraphics[scale=.3]{img/gaddag.png}
\end{figure}
Cette structure est une optimisation du DAWG en terme de
calcul des mots possibles tout en respectant les contraintes
du scrabble. Le GADDAG est donc utilisé dans les algorithmes de
recherche de meilleur coup par IA. Cependant, cette structure prend
plus de place en mémoire que le DAWG.

\subsection{Algorithmes spécifiques}

Nous utilisons deux algorithmes de recherche dans la partie
intelligence artificielle du projet.
\begin{itemize}
    \item \textbf{MiniMax}
    \item \textbf{Monte-Carlo}
\end{itemize}

%----------------------------------------------------------------------------------------------------------------

\section{La liste des besoins visés}

%----------------------------------------------------------------------------------------------------------------

\section{Un agenda prévisionnel des livraisons des besoins}


Pour établir l'agenda prévisionnel, nous allons nous appuyer sur le graphe des dépendances entre les besoins. Ce graphe nous permet d'identifier que certains besoins nécessitent d'autres besoins déjà implémentés. Par exemple, le besoin B01 dépend du besoin B10.

\begin{figure}[H]
    \centering
    \includegraphics[scale=.1]{val/besoin_graph.jpg}
\end{figure}

Nous en concluons donc que nous devons d'abord prioriser l'implémentation des besoins B07, B11 et B12.

Ensuite, nous avons réalisé un sondage auprès des membres de l'équipe afin de déterminer leurs préférences en matière de développement. Certains souhaitent se concentrer sur l'implémentation de l'AI et des algorithmes. Il est également important de noter que certains ne souhaitent pas s'occuper des phases de test, car ils ne sont  pas dans le parcours GL. Il est donc préférable de confier cette tâche aux personnes spécialisées dans ce domaine.

\begin{figure}[H]
    \centering
    \includegraphics[scale=.07]{val/skill.jpg}
\end{figure}

Pour établir notre agenda prévisionnel, nous devons donc combiner le graphe des dépendances avec les préférences de développement des membres de l'équipe.
Mais comment faire ?

\begin{figure}[H]
    \centering
    \includegraphics[scale=.2]{val/black.jpg}
\end{figure}

Nous pourrions réduire ce problème à un problème de coloration de graphe. Cependant, nous avons décidé de ne pas suivre cette approche, car cela nécessiterait trop de temps pour le codage, d'implémentation et les tests. Nous avons estimé que ce développement prendrait au moins une heure, alors que le faire manuellement ne prend que 5 minutes.

\begin{figure}[H]
    \centering
    \includegraphics[scale=.2]{val/graph_coloring.jpg}
\end{figure}

Nous avons donc choisi de combiner ces deux graphes manuellement, en faisant en sorte qu'un membre de l'équipe prenne en charge les dépendances et les associe ensemble.

\begin{figure}[H]
    \centering
    \includegraphics[scale=.2]{val/urss.jpeg}
\end{figure}


Nous obtenons ainsi un graphe qui combine les besoins, leur ordre de développement, les domaines d'intervention dans un projet, ainsi que les membres de l'équipe.

\begin{figure}[H]
    \centering
    \includegraphics[scale=.06]{val/besoin_skill_team.jpg}
\end{figure}


Il ne reste plus qu'une étape pour générer l'agenda prévisionnel: pour chaque besoin, il faut choisir une personne responsable de son développement.

Nous avons donc établi notre agenda prévisionnel.

\begin{figure}[H]
    \centering
    \includegraphics[scale=.3]{val/agenda.png}
\end{figure}

Par exemple, de la Semaine 2 à la Semaine 3, soit 2 semaines de travail, le besoin B07 sera développé par Ludovic et Valentin. Il en va de même pour le besoin B11, qui sera pris en charge par Mathéo et Jean, etc.



%----------------------------------------------------------------------------------------------------------------

\section{Architecture du projet}

\subsection{Modules du projet}

Nous avons choisi une architecture de style MVC pour organiser les
modules du projet. Ce motif d'architecture a été choisi car il 
pose des lignes directrices simples et comprises par tous les membres
vis à vis de la question du placement des différents modules.
De plus, il se prête bien a la mise en place de la partie graphique.
Conformément a ce modèle, nos modules sont donc
découpés en trois parties.\newline
\textit{Modèle}\newline
Le modèle contient les différentes structures de données, 
et autres éléments logiques relatifs aux contraintes du jeu.
Les éléments de cette couche sont: Game, AI.\newline
\textit{Vue}\newline
La vue est constituée de tous les éléments visibles par 
l'utilisateur. Ici, elle est découpée en deux grands axes.
La partie CLI, et la partie GUI.\newline
\textit{Controlleur}\newline
Cette couche de l'architecture traite les interactions avec
l'utilisateur. C'est donc cette partie qui s'occupe d'actualiser
le modèle selon les différentes actions de l'utilisateur.

\subsection{Dépendances du projet}

Le projet est codé en Java version 17.0.13, et est compilé a l'aide du
système de build Maven. Afin de respecter les spécifications du
projet, un certain nombres de dépendances sont ajoutées:
\begin{itemize}
  \item \textbf{JUnit}: Framework de test, utilisé pour les
  différents scénarios de test du code.
  \item \textbf{Jacoco}: Librairie de couverture de code, utilisée
  en conjonction avec JUnit pour certifier le pourcentage de code
  source testé.
  \item \textbf{Commons-cli}: API d'analyse d'options de commandes.
  Utilisé pour interpreter les options passées en ligne de
  commande.
  \item \textbf{JavaFx}: Framework d'interface utilisateur, utilisé
  pour la partie interface graphique.
  \item \textbf{Checkstyle}: Outil d'analyse de code permettant de
  vérifier la conformité du formattage avec le style de codage de 
  Google.
\end{itemize}

%----------------------------------------------------------------------------------------------------------------

\end{document}